\section{Introduction}

Facial affect is commonly estimated as a point in the two-dimensional
valence--arousal (VA) plane of the circumplex model~\cite{russell1980circumplex},
rather than as one of a few discrete categories. A growing line of work attaches
uncertainty to these estimates, through heteroscedastic likelihoods, deep
evidential regression~\cite{amini2020evidential,wu2023deer}, or ensembles, and
converts that uncertainty into prediction intervals by split conformal
calibration~\cite{vovk2005algorithmic,lei2018distribution}. The appeal is a
finite-sample guarantee: on exchangeable data, a $90\%$ interval contains the
truth at least $90\%$ of the time, whatever the model.

That guarantee is only as useful as the event it is stated about. In practice
the intervals are calibrated per axis, one for valence and one for arousal, and
each comes with its own marginal guarantee. But a downstream consumer of a VA
estimate (an abstention rule, a display, a decision) uses the \emph{pair}. The
event that matters is that the true point lies in the predicted region, and two
marginal $90\%$ guarantees imply only $80\%$ for that event. The question of
which event the guarantee is about, and for whom it holds, is the subject of
this paper.

\paragraph{Contributions.}
\begin{enumerate}
  \item \textbf{A diagnosis.} On AffectNet~\cite{mollahosseini2017affectnet},
        per-axis conformal intervals at nominal $90\%$ jointly cover the true VA
        point only $85$--$86\%$ of the time, across every uncertainty method we test.
        Coverage is also uneven across expressions, with the worst seen expression
        at $78$--$84\%$ depending on the predictor.
  \item \textbf{A simple fix with a guarantee.} A single max-normalised
        nonconformity score gives a rectangular VA region with finite-sample
        joint coverage (Proposition~\ref{prop:joint}). It restores nominal joint
        coverage, improves the worst-expression coverage, and is up to $13\%$ smaller
        than the Bonferroni correction that gives the same guarantee.
  \item \textbf{Where the guarantee breaks.} The AffectNet release we study
        contains \emph{contempt} only in its test split, giving a natural unseen
        expression. There the guarantee fails for every method
        ($30$--$39\%$ coverage), and no uncertainty estimator we evaluate
        flags it reliably. This is a concrete, reproducible instance of the
        limits of distribution-free inference~\cite{barber2021limits} in affect.
  \item \textbf{A clean, broader benchmark.} Five uncertainty methods
        (deterministic, Gaussian, deep evidential regression, MC-dropout, deep
        ensembles) over five seeds, scored with proper scoring
        rules~\cite{gneiting2007strictly} as well as calibration error, with a
        conformal calibration split that is never used for model selection.
\end{enumerate}

\paragraph{Scope.} We study still images and a single dataset family. The claims
are about what calibrated VA uncertainty does and does not guarantee, not about
state-of-the-art accuracy, and we position our predictor against published
results in Section~\ref{sec:results}. The system is intended for estimating a
consenting user's own state; Section~\ref{sec:ethics} discusses its limits.
